\section{Preparing for Hurricane Season 2013}

\paragraph{The Ice Caps Are Melting}


Michael Hoffman has described an unusual conspiracy theory that there is a hidden cryptocracy behind seemingly contingent events. They manipulate opinion through symbols. (That is why the correct understanding of symbols is so necessary.) Curiously, they actually reveal what they are doing through popular media, particularly films.

I never took such ideas seriously, but I am now rethinking them. For example, all my life I had assumed that books like 1984 and the Brave New World were warnings to us to avoid such dystopian societies. However, I now think that they were actually preparing us to accept such a world, even willingly. There can be no doubt that we are embracing, in the West, more and more of the features described in that those books.

A creepier video can clinch the deal to any doubters. In 1970, Tiny Tim sang a song to a group of children\footnote{\url{https://www.youtube.com/watch?v=5xfAZ6y9zzM}} about the polar ice caps melting and drowning all the world. (See around 42 or 43 minutes into it). I don't know why the kids did not flip out, but consider that they are now in positions of power. If you believe in the cryptocracy, you will see that the idea of melting ice caps was implanted in to the group consciousness some 40 years ago. If you believe in science, you will see anthropogenic global warming as a great scientific discovery. I'll leave the decision up to you, but take into consideration that Tiny Tim was no scientist.

\paragraph{Politics and Tradition}


Somehow, the idea of Tradition has become intertwined with nominal political movements, at least in some circles. For that to make sense, the goal of such a movement would be the restoration of a traditional society. Given that, we have analyzed to some extent the following works.

\begin{itemize}


\item \textbf{Ananda Coomaraswamy}'s \textit{Spiritual Authority and Temporal Power in Indian Theory of Government}. Although that deals specifically with the Hindu tradition, it serves as a background to other works. We saw, too, how this work corrected Evola's misconception about the relationship between the Brahman caste and the King.

\item \textbf{Rene Guenon}'s \textit{Spiritual Authority and Temporal Power}. This is more abstract than Coomaraswamy's work since it deals with general principles rather than the specific situation in earlier civilizations.

\item \textbf{Guido De Giorgio}'s section on the constitution of a traditional society from \textit{The Roman Tradition}. This is more esoteric, since it describes the inner states associated with the various caste distinctions and the role of the Leader without reference to any specific situations or sacred scriptures. His is the only work that deals effectively with the Worker caste and the role of the Leader.

\item There is \textbf{Julius Evola}'s essay on Monarchy with some important details. For the sake of completeness, we will next make available his essay from the Ur/Krur group with the same title of Spiritual Authority and Temporal Power. I don't believe it has been translated into English, but let me know otherwise.

\item We have mentioned \textbf{Dante}'s essay \textit{On Monarchy}. As a political realist, he realizes that if disputes cannot be resolved rationally or spiritually, then a duel may be the ultimate solution. This applies only in certain defined situations and should be contrasted with Evola's opinion that the result of a duel is a merely contingent event.

\item \textbf{Fustel de Coulanges}' The Ancient City describes the constitution of traditional societies and ancient Greece and Rome.
\end{itemize}


Although there is room for discussion among all these perspectives, there is nevertheless a family resemblance. Primary among them is the recognition of a spiritual unity and none of them begins with racial or ethnic identity. Actually, experience proves this as evidenced by the frustrations expressed by those who start from the bottom up. It is true, by the Law of Affinity, that certain traditions are appropriate to a given peoples, but once the soul and the spirit have been lost, the body loses its coherence, devolving into atomic individuals.

Hence, any political aspirations based on tradition must start with a spiritual impulse. This cannot be defined or delimited in a discourse as its source must needs be transcendent. That is why Guenon emphasized that a transformation to reverse the flow of progressive time is a task possible only for an elite who understand and live by certain principles. De Giorgio goes a step further and points out that there can only be one leader and he represents the spiritual power in the temporal realm. That is why there can ultimately be just one movement and competing movements are part of the counter Tradition.

There is a trend in certain circles to gather disparate material together, so long as they are ``on the right'' in some way, in the vain hope that a unitary movement will ensue. This is the intellectual analog of Mulligan Stew. When Hobos came together, whatever edibles they could gather would be tossed into the pot to make a stew; obviously it would never be the same from day to day. That is food suitable for bums, not for men. Likewise, that applies to ``intellectual'' mulligan stew.

\paragraph{Innocent as Doves}


In a recent episode of Magic City\footnote{\url{http://www.starz.com/originals/magiccity}}, the Cuban trophy wife starts a conversation with her husband, the owner of a Miami Beach hotel.

Wife: In my world, lying and cheating your way to success is not a virtue.

Husband: In my world, I protect those I love.

Wife: Whatever it takes?

Husband: Whatever it takes!

How do they coexist in the same world? One is as innocent as a dove, the other as wise as a serpent. The Leader takes on the sins of the people in order to keep them innocent.

\paragraph{Two Books}


Two books have come to my attention this week; one a history of the movement inaugurated by Guenon, the other a history of Jesus penned by a Muslim. The former is René Guénon: Une politique de l'esprit by David Bisson. I have not seen it, but only partly read a review of it. Unable to conceive of transcendence, the professor reduces Guenon to a metapolitical movement and seeks to categorize him among similar figures of his time.

The premise is incorrect, since Guenon is the latest in a long line of metaphysicians. Tradition, even the Western tradition, is based on the symbolic understanding of scriptures (i.e., beyond the literal interpretation), the perennial philosophy, and a knowingness beyond discursive thought.

That has been largely lost in the West. To beat an old horse, the Protestant Sola Scripture teaching displaced the perennial philosophy. Specifically, the anthropomorphic god gleaned from the literal biblical readings has displaced metaphysics. The Traditional view is that when the literal meaning conflicts with a higher understanding, then a symbolic meaning is more likely. This is beyond the intellectual domain of a professor. Unfortunately, to understand Guenon, a man must embark on the path of spiritual realization.

The other book is Zealot: The Life and Times of Jesus of Nazareth by a Muslim biblical scholar, Reza Aslan. Oddly enough, for an allegedly scholarly book, it has reached \#1 on Amazon. Again, he breaks no new ground, as far as I can tell, but is yet another link in the chain of scholars looking for the true ``historical'' Jesus. The last such book I read several years ago was by a Marcus Borg who envisioned Jesus as a liberal democrat who would have supported universal health care.

Mr. Aslan goes further, asserting that Jesus was some sort of rabble rousing revolutionary figure, willing to overthrow the established order of things. In this, he predated by centuries that spectre that Karl Marx discovered to be haunting Europe. Since everyone who is anyone today sees himself as a zealot and revolutionary, we can well understand the popularity of this book. They are happy to let Jesus onto their team since the holy name still has a certain cachet. Unfortunately, the perpetual revolutionary will eventually run out of things to rebel against, once he becomes aware that he is actually ``the man''. His path can only devolve into chaos as every trace of order is destroyed. For us, on the other hand, Jesus is the incarnation of the Logos, hence the embodiment of the very order of things.


\flrightit{Posted on 2013-07-29 by Cologero}

\begin{center}* * *\end{center}

\begin{footnotesize}
\begin{sffamily}
\texttt{Thomas Blanchard on 2013-07-30 at 10:25 said:}

Forgive my ignorance; what is the source on ``Evola's opinion that the result of a duel is a merely contingent event''? I don't recall this from Revolt, Mystery of the Grail, or any of the articles collected in Metaphysics of War.

\hfill

\texttt{Cologero on 2013-07-30 at 12:31 said:}

Thomas, that is my interpretation of what Evola wrote in Fascism Viewed from the Right, in which he asserted that defeat is a contingent event and could not be related to truth of the principles involved.

Compare to what Dante wrote about the Roman empire. In your research, did you find that Evola claimed that the outcome of a duel reflected divine justice?

\hfill

\texttt{JA on 2013-07-30 at 12:58 said:}

Perhaps there is a difference between a duel between individuals in a traditional society and a war between secular modern states ? Evola did hold in Revolt to the justness of dueling, so the Fascism reference I think refers to a different principle in a different time.

\hfill

\texttt{Jacob on 2013-07-30 at 15:05 said:}

I'm with you, JA.

I understand it as Evola pointing out that a martial victory on the physical plane has no bearing on objective truth. Just because the forces of revolution in France killed the king doesn't mean the king was not rightful ruler of France.

If I understand de Giorgio right (and it's highly possible I do not) the leader is supposed to be as close to pure action as possible. So that means the person most capable in martial respects should be the one in charge right? Then the priestly caste is there to make sure his enforcement of things is correct in terms of Tradition. If so it seems to resolve the question to me.

Is this a correct deduction or did I just read my own biases into de Giorgio's. ideas?

\hfill

\texttt{Izak on 2013-07-30 at 16:38 said:}

Just wanted to say thanks so much for the translations on statecraft!

I'm a bit behind in reading them. I've got some catching up to do.

\hfill

\texttt{Thomas Blanchard on 2013-07-30 at 16:55 said:}

Colgero and JA,

I am specifically referring to this passage from Revolt, in the chapter, ``The Greater and the Lesser Holy War''. The keys here are the last paragraph, and also when he mentions that the action of the trial of strength needs to be given an ``adequate orientation'' (I believe this illustrates the difference you were speaking of, JA.)

------------------------------------

Before continuing, I need to mention an aspect of the traditional spirit that is related to the Law and to the views expounded so far. I am talking about various ordeals of character and so-called divine judgments.

Quite often the test of truth, right, justice, and innocence was made to depend on a trial that consisted of decisive action (experimentum crucis). Just as the law was traditionally believed to have a divine origin, likewise injustice was considered to be a violation of the divine Law and to be detectable through the outcome of a human action that had been given an adequate orientation. A Germanic custom consisted of delving into the divine will through the test of arms as a particular form of oracle mediated by action; the idea that originally was at the basis of the custom of challenging somebody to a duel is not very different. Starting with the principle: ``de coelo est fortitudo'' (Annales Fuldenses), this principle was eventually extended to feuding states and nations. A battle as late as that of Fonteney (A.D. 841) was conceived as a ``divine judgement'' that was invoked to establish the rights of two brothers both claiming the legacy of Charlemagne. When a battle was fought in this spirit, it followed special rules: the winner was forbidden to loot and to exploit strategically and territorially the successful outcome, and both sides were expected to tend equally to the fallen and to the wounded. According to the general view that was preserved through the entire Carolingian period, however, even when the idea of a specific proof was not required, victory and defeat were felt to be signs ``from above'' establishing justice or injustice, truth or guilt.

... 

These views were destined to be regarded as pure superstition wherever ``progress'' systematically deprived the human virtues of any possibility of establishing an objective contact with a superior order of things. Once man's strength was thought to be on the same level as that of animals, that is, as the faculty of mechanical action in a being who is not at all connected to what transcends him as an individual, the trial of strength obviously becomes meaningless and the outcome of every competition becomes entirely contingent and lacking a potential relation with an order of higher ``values.''

------------------------------------

\hfill

\texttt{Michael on 2013-07-30 at 18:38 said:}

Regarding tiptoe through the tulips, I've come to a similar conclusion this year. Things are unfolding in such a way that it is hard to believe that it is not part of a brilliant plan. Who do you think ``they'' are? I personally feel it is just too well done to be a group of people.

\hfill

\texttt{Saladin on 2013-07-30 at 21:51 said:}

Regarding Mr. Aslan: He was born a muslim, then turns evangelical, then becomes muslim again!!!This gentleman ( a compatriot of mine to boot) claims to be a muslim, but he totally ignores the fact that his portrayal of Jesus-christ (which he states is what he believes as a ``historian'') goes against Islamic tradition's view of Jesus which is that of the highest of prophets. Mr. Aslan is neither a

christian nor a Muslim. Given the enormous support he has received from mass media (except Fox News), as well as his other political activities, we can safely guess where his loyalty lies (should I name the tribe?)

Very interesting choice fro a title too. Remember who the ``zealots'' were? (look here for etymology:
\url{http://www.etymonline.com/index.php?term=zealot}). Aslan's claim is that Jesus was a zealot whose ultimate aim was to fight the Roman Empire (and of course help the poor along the way)!! Really?

No, Mr. Aslan is no muslim. He is a clever guy who knows what his masters require of him to be amply rewarded.

\hfill

\texttt{JA on 2013-07-31 at 11:06 said:}

@Thomas that is the quote I was thinking of.

@Jacob -- I think the Capo is the same as Plato's philosopher king, he is an intellectual warrior. In modern society people's developments tend to directed into one area only so we have nerds who have brains but no body strength and we have jocks who are strong but can't add or subtract. I the the Capo as being Caesarian, wise and courageous.

@Saladin and Michael, I used to be an antisemite of the Leon de Poncins or ``Maurice Pinay'' variety but have moved away from that. I've met many Orthodox Jews who share the same values I do and have no wish to partake in worldly affairs. The bad Judaic-derived movements that have so messed the world up (Marxism, Zionism, Freudianism etc) all seem to derive from the Jewish heretic Benedict de Spinoza. Who Spinoza's masters and teachers were I can not say but I think ultimately there is a conspiracy at work over the past several hundred years and it is Satanic literally.

\hfill

\texttt{Pickman on 2013-07-31 at 18:19 said:}

Cologero, not too long ago if I recall, you were very skeptical about the whole symbolism and outward expression of the cryptocracy (or conspiracies and the occult war) through media outlets. Even going so far as deriding those who pointed out these topics. They constantly leave clues everywhere and this is no joke reducible to some ``spooky youtube video''. It wouldn't be a stretch to assume fully that what sci-fi and popular fantasy/horror fiction of today is a prelude in preparation for fact tomorrow (both Wells and Huxleys later works prove this correctly).

I've read Hoffman's erudite analysis on these topics (especially in advertising etc). He only scratch's the surface as only one man can. Other sites have picked up on it, but have not subjected the vast data accumulated into a scholarly treatise aka, Hoffman. This task is laid out for those interested in exposing the tactics of the media cryptocracy or whatever name one wishes to place on the beast.

\hfill

\texttt{Pickman on 2013-07-31 at 18:49 said:}

For those interested in M.Hoffman's ``Secret Societies and Psychological Warfare'' it can be viewed here: downloads.umu.nu/Books/Michael.A.Hoffman.II--Secret.Societies.And.Psychological.Warfare. 1992 .pdf , I'd also recommend buying it on paperback to support the author.

``Who then is the modern man? He is a mind-bombed patsy who gets his marching orders from ``twilight language'' key words sprinkled throughout ``his'' news and current events. Even as he dances to the tune of the elite managers of human behavior, he scoffs with great derision at the idea of the existence and operation of a technology of mass mind control emanating from the media and government. Modern man is the ideal hypnotic subject: puffed up on the idea that he is the crown of creation, he vehemently denies the power of the hypnotist's control over him, even as his head bobs up and down on a string.'' -- M.H.

Word's of wisdom.

\hfill

\texttt{Ash on 2013-08-01 at 18:57 said:}

This article reminded me on your comments about an anti-Traditional elite. It appears that the new anti-porn laws in the UK will be blocking other things too, like websites depicting violence, smoking, and what I found most interesting, ``esoteric materials''.

Now I'm a bit doubtful as to whether whoever inserted that phrase was specifically thinking of the works of Guenon (I would imagine that they meant in the sense of ``hidden''), but I find it rather disconcerting that this was specifically included, particularly with the lack of definition. I can see how the others might fall in under the nanny state's gaze, but to block access to ``esoteric materials''? Perhaps something to keep an eye on. It is harder to fight a hidden enemy than an open one, and to do battle when you have forgotten there is an enemy at all is nigh impossible.

\url{http://www.huffingtonpost.com/2013/07/29/uk-internet-filter-block-more-than-porn\_n\_3670771.html?utm\_hp\_ref=world}


\hfill

\texttt{White Rabbit on 2015-07-12 at 07:54 said:}

I think you're being unfair to the bums. We can probably assume they don't imagine themselves to be fine chefs.


\hfill

\end{sffamily}
\end{footnotesize}

